\documentclass[12pt,a4paper]{article}

\usepackage[utf8]{inputenc}
\usepackage[T1]{fontenc}
\usepackage[margin=2.5cm]{geometry}
\usepackage{booktabs}
\usepackage{hyperref}
\usepackage{parskip}
\usepackage{titlesec}
\usepackage{enumitem}
\usepackage{listings}
\usepackage{xcolor}
\usepackage{array}

\hypersetup{
    colorlinks=true,
    linkcolor=blue,
    urlcolor=blue,
    citecolor=blue
}

\lstset{
    basicstyle=\ttfamily\small,
    backgroundcolor=\color{gray!10},
    frame=single,
    breaklines=true,
    columns=fullflexible
}

\title{
    \textbf{CENG 467 — Term Project Progress Report} \\[0.5em]
    \Large Advanced RAG with Hard Negative Mining
}

\author{
    Erman Akıncı (300201066) \quad
    Ahmet Yılmaz (300201067) \quad
    Kaan Cesur (290201066)
}

\date{
    İzmir Institute of Technology (İYTE) \\
    CENG 467 Natural Language Understanding and Generation \\
    Prof.\ Dr.\ Aytuğ ONAN \quad Spring 2026 \\[0.5em]
    \textit{Submitted: May 14, 2026}
}

\begin{document}

\maketitle
\tableofcontents
\newpage

% ---------------------------------------------------------------
\section{Project Title, Group Members, and Problem Definition}

\subsection*{Task Description}

This project addresses the task of \textbf{open-domain question answering (QA)} using a
Retrieval-Augmented Generation (RAG) pipeline. Given a natural language question, the
system must retrieve relevant passages from a corpus and produce a correct short answer.

\subsection*{Why This Task Is Important}

Large language models (LLMs) are prone to hallucination when answering factual questions
from parametric memory alone. RAG systems ground model outputs in retrieved evidence,
improving factual accuracy and interpretability. The retrieval component is a critical
bottleneck: a retriever that fails to surface the correct passage makes correct answer
generation impossible regardless of generator quality.

\subsection*{Input and Output}

\begin{itemize}[noitemsep]
    \item \textbf{Input:} A natural language question
          (e.g., \textit{``Were Scott Derrickson and Ed Wood of the same nationality?''})
    \item \textbf{Output:} A short answer string (e.g., \textit{``yes''})
    \item \textbf{Supporting context:} A set of passages retrieved from a candidate corpus
\end{itemize}

\subsection*{Task Type}

Multi-hop question answering with retrieval. The task requires combining information
from multiple documents, making it a joint retrieval and reading comprehension problem.

\subsection*{Technical Challenge}

Standard retrieval methods such as BM25 rely on lexical overlap, which fails for
paraphrased or semantically complex questions. Dense retrievers learn semantic
representations but require hard negative examples during training to distinguish
between superficially similar but factually incorrect passages. Without hard negative
mining, the retriever learns trivially from easy negatives and degrades on difficult
multi-hop questions.

% ---------------------------------------------------------------
\section{Dataset / Benchmark Status}

\subsection*{Dataset Information}

\begin{table}[h]
\centering
\begin{tabular}{ll}
\toprule
\textbf{Property} & \textbf{Value} \\
\midrule
Name              & HotpotQA \\
Setting           & Distractor \\
Source            & HuggingFace \texttt{datasets} library \\
Train split       & 90,447 examples \\
Validation split  & 7,405 examples \\
Test split        & Not publicly available \\
Input format      & Question + 10 candidate passages (title + sentences) \\
Output format     & Short answer string \\
Answer types      & Span extraction, yes/no \\
\bottomrule
\end{tabular}
\caption{HotpotQA dataset properties}
\end{table}

\subsection*{Current Status}

The dataset has been successfully downloaded and loaded via the HuggingFace
\texttt{datasets} library. Preprocessing and tokenization for BM25 retrieval is complete.
Context encoder tokenization for DPR (max length 256 tokens, padding, truncation) is
implemented.

\subsection*{Preprocessing Steps Completed}

\begin{itemize}[noitemsep]
    \item Passage construction: document title concatenated with all sentences
    \item BM25 tokenization: lowercasing, punctuation removal, whitespace split
    \item DPR tokenization: HuggingFace tokenizer, truncation at 256 tokens for contexts
          and 64 tokens for questions
\end{itemize}

\subsection*{Known Issues}

\begin{itemize}[noitemsep]
    \item Test set labels are not publicly released; all evaluation uses the validation split.
    \item Some passages exceed 256 tokens and are truncated during DPR encoding.
    \item The distractor setting provides only 10 passages per question rather than a
          full corpus; scaling to full corpus retrieval is a planned extension.
\end{itemize}

% ---------------------------------------------------------------
\section{Literature Review Progress}

\paragraph{Dense Passage Retrieval (DPR) — Karpukhin et al., 2020.}
DPR introduced bi-encoder retrieval where questions and passages are encoded
independently using BERT-based encoders and matched via dot product similarity.
This work established the foundation for our dense retrieval baseline and demonstrated
that dense retrievers can outperform BM25 when trained with in-batch negatives.

\paragraph{ANCE — Xiong et al., 2021.}
Approximate Nearest Neighbor Negative Contrastive Estimation (ANCE) trains dense
retrievers with hard negatives selected from the top results of the retriever itself
at each training checkpoint. This approach is the primary methodological reference
for the hard negative mining component of our proposed system.

\paragraph{RocketQA — Qu et al., 2021.}
RocketQA improved DPR training by using cross-encoder-based hard negative filtering
and denoised training data. It provides a reference upper bound for retrieval
performance and informs our ablation study design.

\paragraph{HotpotQA — Yang et al., 2018.}
The original paper introducing the dataset describes the multi-hop reasoning challenge
and the distractor setting, and reports baseline numbers for BM25 and early neural
retrieval systems.

% ⚠️ BOŞLUK: Erman — 1-2 makale daha ekle (örn. RAG Lewis et al. 2020, ColBERT)

\subsection*{Technical Direction}

Based on the reviewed literature, our project follows the DPR + hard negative mining
paradigm. We will use BM25-retrieved passages as hard negatives for fine-tuning the
DPR model, following the finding in ANCE that BM25 negatives are a strong and
computationally inexpensive source of hard negatives.

% ---------------------------------------------------------------
\section{Baseline Models or Prompting Strategies}

\subsection*{Baseline 1: BM25}

BM25 (Okapi BM25) is a classical probabilistic retrieval model based on term frequency
and inverse document frequency. It requires no training and serves as a strong lexical
baseline. We use the \texttt{rank\_bm25} Python library with lowercased,
punctuation-removed tokenization.

\textbf{Status: Fully implemented and evaluated.}

\subsection*{Baseline 2: Vanilla DPR}

We use the pretrained DPR encoders released by Facebook Research
(\texttt{facebook/dpr-question\_encoder-single-nq-base} and
\texttt{facebook/dpr-ctx\_encoder-single-nq-base}), trained on Natural Questions.
No fine-tuning is applied. This baseline measures off-the-shelf dense retrieval
performance on HotpotQA without domain adaptation or hard negative training.

\textbf{Status: Implemented, evaluation in progress.}

\subsection*{Why These Baselines}

BM25 represents the lexical retrieval paradigm and is a standard baseline in all
retrieval-augmented QA literature. Vanilla DPR represents the semantic retrieval
paradigm without hard negative fine-tuning, allowing us to isolate the contribution
of hard negative mining in the final system.

% ---------------------------------------------------------------
\section{Initial Experimental Results}

Table~\ref{tab:results} reports retrieval results on 500 examples from the
HotpotQA validation split.

\begin{table}[h]
\centering
\begin{tabular}{llllp{3.5cm}}
\toprule
\textbf{Model} & \textbf{Recall@5} & \textbf{F1} & \textbf{Exact Match} & \textbf{Status} \\
\midrule
BM25                    & 0.972 & 0.060 & 0.004 & Complete \\
DPR (vanilla)           & 0.982 & 0.064 & 0.002 & Complete \\
DPR + Hard Negatives    & —     & —     & —     & Planned \\
\bottomrule
\end{tabular}
\caption{Retrieval results on HotpotQA validation (500 examples, top-k=5)}
\label{tab:results}
\end{table}


\subsection*{Interpretation}

BM25 achieves a Recall@5 of 0.972, meaning it successfully retrieves at least one
gold passage in the top 5 results for 97.2\% of questions. The low F1 (0.060) and
Exact Match (0.004) scores reflect the absence of a reader component: the predicted
answer is approximated as the first sentence of the top retrieved passage, which is
rarely a short exact-match string. These scores will improve once a reader model or
LLM is integrated in the next phase.

% ---------------------------------------------------------------
\section{Planned Improvements and Technical Direction}

The core contribution of this project is training a DPR model with hard negative mining.
The planned pipeline proceeds as follows:

\begin{enumerate}[noitemsep]
    \item \textbf{Hard Negative Generation:} For each training question, run BM25 to
          retrieve top-$k$ passages. Passages that score highly under BM25 but do not
          contain the gold answer are selected as hard negatives.
    \item \textbf{DPR Fine-tuning:} Fine-tune the DPR question and context encoders
          using contrastive loss (in-batch negatives + hard negatives) on the HotpotQA
          training split (90k examples).
    \item \textbf{Reader Integration:} Pass the top-3 retrieved passages together with
          the question to an LLM or extractive QA model to generate the final answer.
    \item \textbf{FAISS Index:} Build a FAISS flat index over all retrieved passages for
          scalable nearest-neighbor search.
\end{enumerate}

% ---------------------------------------------------------------
\section{Ablation or Prompt Sensitivity Plan}

We plan the following ablations to isolate the contribution of each component:

\begin{table}[h]
\centering
\begin{tabular}{ll}
\toprule
\textbf{Ablation} & \textbf{Variable} \\
\midrule
A1 & BM25 vs.\ DPR (no fine-tuning): lexical vs.\ semantic retrieval \\
A2 & DPR vanilla vs.\ DPR + hard negatives: effect of mining \\
A3 & Hard negative count: $k \in \{1, 5, 10\}$ per question \\
A4 & Top-$k$ retrieval: Recall@1 vs.\ Recall@3 vs.\ Recall@5 \\
A5 & With reader vs.\ without reader: retrieval-only vs.\ full RAG \\
\bottomrule
\end{tabular}
\caption{Planned ablation study}
\end{table}

% ---------------------------------------------------------------
\section{Error Analysis and Generation Quality Analysis}

\subsection*{Planned Error Analysis}

We will sample 50 failure cases from each baseline and categorize them as follows:

\begin{itemize}[noitemsep]
    \item \textbf{Lexical gap failures:} BM25 fails because the question uses different
          vocabulary than the passage.
    \item \textbf{Semantic confusion:} DPR retrieves topically related but factually
          incorrect passages.
    \item \textbf{Multi-hop failures:} Neither baseline retrieves both required supporting
          passages simultaneously.
    \item \textbf{Yes/no failures:} Binary-answer questions where the retrieved passage
          does not contain an explicit yes/no signal.
\end{itemize}

\subsection*{Generation Quality Plan}

Once the reader is integrated, we will evaluate:
factual correctness relative to gold answers, fluency of generated answers, and
hallucination rate (answers not grounded in retrieved passages).

% ---------------------------------------------------------------
\section{Ethical Considerations and Bias}

This project uses the HotpotQA benchmark derived from Wikipedia. Potential concerns
include:

\begin{itemize}[noitemsep]
    \item \textbf{Hallucination risk:} LLM-based readers may generate plausible but
          factually incorrect answers not supported by retrieved passages.
    \item \textbf{Coverage bias:} Wikipedia overrepresents English-language and
          Western-centric topics, which may affect retrieval performance on questions
          about underrepresented subjects.
    \item \textbf{Evaluation limitations:} Exact Match and F1 against a single reference
          answer may penalize semantically correct paraphrases.
    \item \textbf{API cost constraints:} If an external LLM API is used for answer
          generation, cost limits the number of evaluated examples.
\end{itemize}

% ---------------------------------------------------------------
\section{GitHub and Reproducibility Status}

\textbf{Repository:} % ⚠️ BOŞLUK — GitHub linkini buraya ekle
\url{https://github.com/KULLANICI_ADI/CENG467_Term_Project}

\subsection*{Current Repository Structure}

\begin{lstlisting}
CENG467_Term_Project/
├── src/
│   ├── load_dataset.py       # HotpotQA loading and inspection
│   ├── baseline_bm25.py      # BM25 retrieval baseline
│   └── baseline_dpr.py       # DPR retrieval baseline
├── results/
│   └── bm25_results.json     # BM25 evaluation output
├── data/                     # Empty (loaded via HuggingFace)
├── notebooks/                # Planned
├── requirements.txt
├── .gitignore
└── README.md
\end{lstlisting}

\subsection*{Reproducibility Instructions}

\begin{lstlisting}[language=bash]
# Install dependencies
pip install -r requirements.txt

# Run BM25 baseline (downloads dataset automatically)
python src/baseline_bm25.py

# Run DPR baseline (downloads model weights, ~800 MB)
python src/baseline_dpr.py
\end{lstlisting}

\subsection*{Status Checklist}

\begin{table}[h]
\centering
\begin{tabular}{ll}
\toprule
\textbf{Item} & \textbf{Status} \\
\midrule
GitHub repository       & % ⚠️ BOŞLUK
                          Open / link pending \\
\texttt{requirements.txt} & Done \\
README                  & Done \\
Dataset loading script  & Done \\
BM25 baseline script    & Done \\
DPR baseline script     & Done \\
\texttt{results/bm25\_results.json} & Done \\
Consistent commit history & In progress \\
\bottomrule
\end{tabular}
\caption{Repository status checklist}
\end{table}

% ---------------------------------------------------------------
\section{Current Challenges and Next Steps}

\subsection*{Current Challenges}

\begin{itemize}[noitemsep]
    \item \textbf{Computational cost:} DPR passage encoding is slow on CPU
          (approximately 5--10 minutes for 500 examples). Fine-tuning will require
          GPU access (Google Colab or university cluster).
    \item \textbf{No reader component:} Current F1/EM scores reflect retrieval-only
          heuristics, not true QA performance.
    \item \textbf{Hard negative mining not yet implemented:} This is the core
          contribution and the focus of the next development phase.
\end{itemize}

\subsection*{Next Steps Before Final Submission}

\begin{enumerate}[noitemsep]
    \item Complete DPR vanilla evaluation and record results in Table~\ref{tab:results}
    \item Implement hard negative mining using BM25 scores on the training split
    \item Fine-tune DPR with hard negatives on GPU
    \item Integrate a reader model (extractive QA or LLM API)
    \item Run the full ablation study
    \item Write the final report in LNCS format
\end{enumerate}

\end{document}
